\section{Artifact Design and
Implementation}\label{artifact-design-and-implementation}

The implementation uses a canonical event-log schema with five fields:
case identifier, activity, resource, start timestamp, and end timestamp.
Table~\ref{tab:schemas} summarizes the three output schemas used by the artifact. The
simulated event log is kept minimal so that it remains compatible with
conventional process-mining and BPS evaluation tools.
The reasoning and handover logs are additional outputs of the LLM-agent
proxy condition and are not required by the baselines.

\begin{longtable}[]{@{}
  >{\raggedright\arraybackslash}p{(\columnwidth - 4\tabcolsep) * \real{0.3333}}
  >{\raggedright\arraybackslash}p{(\columnwidth - 4\tabcolsep) * \real{0.3333}}
  >{\raggedright\arraybackslash}p{(\columnwidth - 4\tabcolsep) * \real{0.3333}}@{}}
\caption{Output schemas produced by the simulator.}
\label{tab:schemas}\\
\toprule\noalign{}
\begin{minipage}[b]{\linewidth}\raggedright
Output
\end{minipage} & \begin{minipage}[b]{\linewidth}\raggedright
Fields
\end{minipage} & \begin{minipage}[b]{\linewidth}\raggedright
Purpose
\end{minipage} \\
\midrule\noalign{}
\endhead
\bottomrule\noalign{}
\endlastfoot
Simulated event log (L') & \texttt{case\_id}, \texttt{activity},
\texttt{resource}, \texttt{start\_time}, \texttt{end\_time} &
Process-mining-compatible generated log used for quantitative
evaluation. \\
Reasoning log (R') & \texttt{case\_id}, \texttt{step\_index},
\texttt{mode}, \texttt{previous\_resource}, \texttt{activity},
\texttt{selected\_resource}, \texttt{feasible\_resource\_count},
\texttt{feasible\_resources\_top}, \texttt{historical\_prior},
\texttt{selection\_rule}, \texttt{reason} & Local decision trace emitted
by the LLM-agent proxy. \\
Handover log (H') & \texttt{case\_id}, \texttt{from\_resource},
\texttt{to\_resource}, \texttt{activity}, \texttt{timestamp},
\texttt{message} & Cross-resource transition trace emitted when a case
moves between resources. \\
\end{longtable}

From the training split, the simulator learns activity-resource
capabilities, service-time samples per activity/resource pair, empirical
inter-event waiting-time samples, trace variant frequencies, case
inter-arrival samples, and handover frequencies between resources. These
learned components are stored as a profile object and shared by the
simulation conditions so that policy differences can be compared under a
common data foundation.

The simulation procedure has four stages. First, the input log is
normalized to the canonical schema and split into training and test
logs. Second, the training log is transformed into log-derived profiles:
resource capabilities, duration samples, waiting-time samples,
case-arrival samples, trace variant counts, activity counts, resource
counts, and handover priors. Third, each policy simulates the same
number of cases as the held-out test log, beginning at the first
timestamp of the test window and sampling empirical case inter-arrival
times. Fourth, the generated logs are evaluated against the held-out
test log with lightweight metrics and the Chapela-Campa distance
measures.

The three main simulation conditions differ only in the resource
decision policy. This isolates the comparison: if the timing model or
trace sampler differed across conditions, performance differences could
not be attributed to the agent decision layer. Table~\ref{tab:conditions}
summarizes the policy comparison.

\begin{longtable}[]{@{}
  >{\raggedright\arraybackslash}p{(\columnwidth - 6\tabcolsep) * \real{0.2500}}
  >{\raggedright\arraybackslash}p{(\columnwidth - 6\tabcolsep) * \real{0.2500}}
  >{\raggedright\arraybackslash}p{(\columnwidth - 6\tabcolsep) * \real{0.2500}}
  >{\raggedright\arraybackslash}p{(\columnwidth - 6\tabcolsep) * \real{0.2500}}@{}}
\caption{Simulation conditions and resource decision policies.}
\label{tab:conditions}\\
\toprule\noalign{}
\begin{minipage}[b]{\linewidth}\raggedright
Condition
\end{minipage} & \begin{minipage}[b]{\linewidth}\raggedright
Decision locus
\end{minipage} & \begin{minipage}[b]{\linewidth}\raggedright
Resource selection rule
\end{minipage} & \begin{minipage}[b]{\linewidth}\raggedright
Additional outputs
\end{minipage} \\
\midrule\noalign{}
\endhead
\bottomrule\noalign{}
\endlastfoot
\texttt{central\_baseline} & Centralized scheduler & Uniformly samples
from resources historically capable of the activity. & Simulated event
log only. \\
\texttt{agent\_profile} & Resource-agent profile & Samples from
historically capable resources weighted by activity-resource frequency.
& Simulated event log only. \\
\texttt{llm\_agent\_proxy} & Constrained local LLM-style policy & Uses
handover priors when available, otherwise activity-resource priors;
samples among feasible resources with an overuse penalty relative to
historical target shares. & Simulated event log, reasoning log, handover
log. \\
\texttt{llm\_agent\_real} & Optional API-backed LLM policy & Sends the
same constrained local state and feasible resource set to an
OpenAI-compatible chat-completions endpoint; validates JSON output and
falls back to the guarded proxy if the call fails or returns an
infeasible resource. & Simulated event log, reasoning log, handover log,
LLM call and fallback diagnostics. \\
\end{longtable}

The \texttt{llm\_agent\_proxy} can influence local resource allocation
and handover behaviour, but it cannot create activities, resources, or
timestamps outside the learned event-log structure. The optional
\texttt{llm\_agent\_real} condition uses the same constrained state and
feasible resource set. It requires a JSON resource choice and rationale
from an OpenAI-compatible endpoint, and falls back to the guarded proxy
if the call fails or returns an infeasible resource. The reported
experiments use the proxy so that the quantitative comparison remains
reproducible.

The implementation was refined through validation-driven iterations. An
early real-log smoke test generated plausible control-flow metrics but
poor temporal fidelity because simulated cases were too short. The next
version added empirical inter-event waiting times, reducing mean
cycle-time relative error from approximately 98.6\% to approximately
4.2\% for non-LLM baselines on AcademicCredentials. A later evaluation
identified a second timing issue: generated logs were placed in a 2026
simulation window whereas the held-out test log started in 2016, making
absolute event and case-arrival distances uninformative. The simulator
was corrected to start at the first held-out timestamp and to sample
empirical case inter-arrivals from the training log. These iterations
illustrate the role of validation as part of artifact construction.
